% Options for packages loaded elsewhere
% Options for packages loaded elsewhere
\PassOptionsToPackage{unicode}{hyperref}
\PassOptionsToPackage{hyphens}{url}
\PassOptionsToPackage{dvipsnames,svgnames,x11names}{xcolor}
%
\documentclass[
  spanish,
  letterpaper,
  DIV=11,
  numbers=noendperiod]{scrartcl}
\usepackage{xcolor}
\usepackage{amsmath,amssymb}
\setcounter{secnumdepth}{5}
\usepackage{iftex}
\ifPDFTeX
  \usepackage[T1]{fontenc}
  \usepackage[utf8]{inputenc}
  \usepackage{textcomp} % provide euro and other symbols
\else % if luatex or xetex
  \usepackage{unicode-math} % this also loads fontspec
  \defaultfontfeatures{Scale=MatchLowercase}
  \defaultfontfeatures[\rmfamily]{Ligatures=TeX,Scale=1}
\fi
\usepackage{lmodern}
\ifPDFTeX\else
  % xetex/luatex font selection
\fi
% Use upquote if available, for straight quotes in verbatim environments
\IfFileExists{upquote.sty}{\usepackage{upquote}}{}
\IfFileExists{microtype.sty}{% use microtype if available
  \usepackage[]{microtype}
  \UseMicrotypeSet[protrusion]{basicmath} % disable protrusion for tt fonts
}{}
\makeatletter
\@ifundefined{KOMAClassName}{% if non-KOMA class
  \IfFileExists{parskip.sty}{%
    \usepackage{parskip}
  }{% else
    \setlength{\parindent}{0pt}
    \setlength{\parskip}{6pt plus 2pt minus 1pt}}
}{% if KOMA class
  \KOMAoptions{parskip=half}}
\makeatother
% Make \paragraph and \subparagraph free-standing
\makeatletter
\ifx\paragraph\undefined\else
  \let\oldparagraph\paragraph
  \renewcommand{\paragraph}{
    \@ifstar
      \xxxParagraphStar
      \xxxParagraphNoStar
  }
  \newcommand{\xxxParagraphStar}[1]{\oldparagraph*{#1}\mbox{}}
  \newcommand{\xxxParagraphNoStar}[1]{\oldparagraph{#1}\mbox{}}
\fi
\ifx\subparagraph\undefined\else
  \let\oldsubparagraph\subparagraph
  \renewcommand{\subparagraph}{
    \@ifstar
      \xxxSubParagraphStar
      \xxxSubParagraphNoStar
  }
  \newcommand{\xxxSubParagraphStar}[1]{\oldsubparagraph*{#1}\mbox{}}
  \newcommand{\xxxSubParagraphNoStar}[1]{\oldsubparagraph{#1}\mbox{}}
\fi
\makeatother


\usepackage{longtable,booktabs,array}
\usepackage{calc} % for calculating minipage widths
% Correct order of tables after \paragraph or \subparagraph
\usepackage{etoolbox}
\makeatletter
\patchcmd\longtable{\par}{\if@noskipsec\mbox{}\fi\par}{}{}
\makeatother
% Allow footnotes in longtable head/foot
\IfFileExists{footnotehyper.sty}{\usepackage{footnotehyper}}{\usepackage{footnote}}
\makesavenoteenv{longtable}
\usepackage{graphicx}
\makeatletter
\newsavebox\pandoc@box
\newcommand*\pandocbounded[1]{% scales image to fit in text height/width
  \sbox\pandoc@box{#1}%
  \Gscale@div\@tempa{\textheight}{\dimexpr\ht\pandoc@box+\dp\pandoc@box\relax}%
  \Gscale@div\@tempb{\linewidth}{\wd\pandoc@box}%
  \ifdim\@tempb\p@<\@tempa\p@\let\@tempa\@tempb\fi% select the smaller of both
  \ifdim\@tempa\p@<\p@\scalebox{\@tempa}{\usebox\pandoc@box}%
  \else\usebox{\pandoc@box}%
  \fi%
}
% Set default figure placement to htbp
\def\fps@figure{htbp}
\makeatother



\ifLuaTeX
\usepackage[bidi=basic]{babel}
\else
\usepackage[bidi=default]{babel}
\fi
% get rid of language-specific shorthands (see #6817):
\let\LanguageShortHands\languageshorthands
\def\languageshorthands#1{}


\setlength{\emergencystretch}{3em} % prevent overfull lines

\providecommand{\tightlist}{%
  \setlength{\itemsep}{0pt}\setlength{\parskip}{0pt}}



 


\KOMAoption{captions}{tableheading}
\makeatletter
\@ifpackageloaded{tcolorbox}{}{\usepackage[skins,breakable]{tcolorbox}}
\@ifpackageloaded{fontawesome5}{}{\usepackage{fontawesome5}}
\definecolor{quarto-callout-color}{HTML}{909090}
\definecolor{quarto-callout-note-color}{HTML}{0758E5}
\definecolor{quarto-callout-important-color}{HTML}{CC1914}
\definecolor{quarto-callout-warning-color}{HTML}{EB9113}
\definecolor{quarto-callout-tip-color}{HTML}{00A047}
\definecolor{quarto-callout-caution-color}{HTML}{FC5300}
\definecolor{quarto-callout-color-frame}{HTML}{acacac}
\definecolor{quarto-callout-note-color-frame}{HTML}{4582ec}
\definecolor{quarto-callout-important-color-frame}{HTML}{d9534f}
\definecolor{quarto-callout-warning-color-frame}{HTML}{f0ad4e}
\definecolor{quarto-callout-tip-color-frame}{HTML}{02b875}
\definecolor{quarto-callout-caution-color-frame}{HTML}{fd7e14}
\makeatother
\makeatletter
\@ifpackageloaded{caption}{}{\usepackage{caption}}
\AtBeginDocument{%
\ifdefined\contentsname
  \renewcommand*\contentsname{Tabla de contenidos}
\else
  \newcommand\contentsname{Tabla de contenidos}
\fi
\ifdefined\listfigurename
  \renewcommand*\listfigurename{Listado de Figuras}
\else
  \newcommand\listfigurename{Listado de Figuras}
\fi
\ifdefined\listtablename
  \renewcommand*\listtablename{Listado de Tablas}
\else
  \newcommand\listtablename{Listado de Tablas}
\fi
\ifdefined\figurename
  \renewcommand*\figurename{Figura}
\else
  \newcommand\figurename{Figura}
\fi
\ifdefined\tablename
  \renewcommand*\tablename{Tabla}
\else
  \newcommand\tablename{Tabla}
\fi
}
\@ifpackageloaded{float}{}{\usepackage{float}}
\floatstyle{ruled}
\@ifundefined{c@chapter}{\newfloat{codelisting}{h}{lop}}{\newfloat{codelisting}{h}{lop}[chapter]}
\floatname{codelisting}{Listado}
\newcommand*\listoflistings{\listof{codelisting}{Listado de Listados}}
\makeatother
\makeatletter
\makeatother
\makeatletter
\@ifpackageloaded{caption}{}{\usepackage{caption}}
\@ifpackageloaded{subcaption}{}{\usepackage{subcaption}}
\makeatother
\usepackage{bookmark}
\IfFileExists{xurl.sty}{\usepackage{xurl}}{} % add URL line breaks if available
\urlstyle{same}
\hypersetup{
  pdftitle={Caso 07: Inspecciones sanitarias de establecimientos en Chicago},
  pdfauthor={Graciela Gómez; Benjamín Gutiérrez; Juan Pablo Mondaca; Benjamín Vargas},
  pdflang={es},
  colorlinks=true,
  linkcolor={blue},
  filecolor={Maroon},
  citecolor={Blue},
  urlcolor={Blue},
  pdfcreator={LaTeX via pandoc}}


\title{Caso 07: Inspecciones sanitarias de establecimientos en Chicago}
\usepackage{etoolbox}
\makeatletter
\providecommand{\subtitle}[1]{% add subtitle to \maketitle
  \apptocmd{\@title}{\par {\large #1 \par}}{}{}
}
\makeatother
\subtitle{EYP3407 Tópicos Aplicados en Estadística · Trabajo final}
\author{Graciela Gómez \and Benjamín Gutiérrez \and Juan Pablo
Mondaca \and Benjamín Vargas}
\date{2026-11-09}
\begin{document}
\maketitle


\section{Resumen ejecutivo}\label{sec-resumen}

\begin{tcolorbox}[enhanced jigsaw, toptitle=1mm, bottomrule=.15mm, breakable, opacityback=0, title=\textcolor{quarto-callout-important-color}{\faExclamation}\hspace{0.5em}{Qué exige}, toprule=.15mm, bottomtitle=1mm, colframe=quarto-callout-important-color-frame, rightrule=.15mm, titlerule=0mm, colback=white, arc=.35mm, left=2mm, opacitybacktitle=0.6, coltitle=black, leftrule=.75mm, colbacktitle=quarto-callout-important-color!10!white]

Máximo 250 palabras: problema, resultado principal, recomendación y
principal limitación. Se escribe al final.

\end{tcolorbox}

\begin{tcolorbox}[enhanced jigsaw, toptitle=1mm, bottomrule=.15mm, breakable, opacityback=0, title=\textcolor{quarto-callout-note-color}{\faInfo}\hspace{0.5em}{Notas del equipo}, toprule=.15mm, bottomtitle=1mm, colframe=quarto-callout-note-color-frame, rightrule=.15mm, titlerule=0mm, colback=white, arc=.35mm, left=2mm, opacitybacktitle=0.6, coltitle=black, leftrule=.75mm, colbacktitle=quarto-callout-note-color!10!white]

\begin{itemize}
\tightlist
\item
\end{itemize}

\end{tcolorbox}

\section{Problema y datos}\label{sec-datos}

\subsection{Formulación del problema}\label{formulaciuxf3n-del-problema}

Entre los establecimientos de comida de Chicago que reciben una
inspección rutinaria, ¿cuáles tienen mayor probabilidad de reprobar? El
análisis está dirigido al Departamento de Salud Pública de Chicago
(CDPH) y podría apoyar la decisión de qué locales visitar primero al
programar estas inspecciones.

La población son las inspecciones rutinarias de tipo \emph{Canvass}
realizadas entre 2019 y 2024. La unidad de observación es una
inspección: un mismo local aporta tantas observaciones como inspecciones
\emph{Canvass} haya recibido. La respuesta es binaria y vale 1 si el
resultado es \emph{Fail} y 0 si es \emph{Pass} o \emph{Pass w/
Conditions}; el modelo entrega, para cada inspección, una probabilidad
de \emph{Fail} entre 0 y 1.

La predicción se emite el día anterior a la visita, cuando la inspección
ya está programada, y solo usa información disponible en ese momento:
características del local, su ubicación agregada, la fecha programada y
los antecedentes de visitas anteriores. Se trata, por lo tanto, de
predecir un resultado futuro y no de caracterizar retrospectivamente
inspecciones ya realizadas.

Los datos solo incluyen locales inspeccionados, y la frecuencia de
visita depende de su nivel de riesgo. Las conclusiones se aplican a esta
población y no estiman la seguridad sanitaria de todos los locales de la
ciudad. Las asociaciones encontradas son predictivas y no permiten
afirmar que una característica cause el resultado.

\subsection{Procedencia}\label{procedencia}

\begin{tcolorbox}[enhanced jigsaw, toptitle=1mm, bottomrule=.15mm, breakable, opacityback=0, title=\textcolor{quarto-callout-important-color}{\faExclamation}\hspace{0.5em}{Qué exige (enunciado 2.2.1)}, toprule=.15mm, bottomtitle=1mm, colframe=quarto-callout-important-color-frame, rightrule=.15mm, titlerule=0mm, colback=white, arc=.35mm, left=2mm, opacitybacktitle=0.6, coltitle=black, leftrule=.75mm, colbacktitle=quarto-callout-important-color!10!white]

Institución responsable, versión o fecha de descarga, archivos
originales, filtros aplicados, condiciones de uso y diccionario de
variables. Los originales se conservan sin editar.

\end{tcolorbox}

Los datos provienen del conjunto \emph{Food Inspections} del portal de
datos abiertos de la ciudad de Chicago, producido por el Chicago
Department of Public Health. Se utiliza la copia preparada por el curso
el 12 de octubre de 2026, que contiene 104.468 inspecciones realizadas
entre el 02-01-2019 y el 31-12-2024.

\begin{tcolorbox}[enhanced jigsaw, toptitle=1mm, bottomrule=.15mm, breakable, opacityback=0, title=\textcolor{quarto-callout-note-color}{\faInfo}\hspace{0.5em}{Notas del equipo}, toprule=.15mm, bottomtitle=1mm, colframe=quarto-callout-note-color-frame, rightrule=.15mm, titlerule=0mm, colback=white, arc=.35mm, left=2mm, opacitybacktitle=0.6, coltitle=black, leftrule=.75mm, colbacktitle=quarto-callout-note-color!10!white]

\begin{itemize}
\tightlist
\item
  Pendiente: revisar las condiciones de reutilización del portal y citar
  la fuente con formato de referencia.
\end{itemize}

\end{tcolorbox}

\subsection{Construcción de la base
analítica}\label{construcciuxf3n-de-la-base-analuxedtica}

\begin{tcolorbox}[enhanced jigsaw, toptitle=1mm, bottomrule=.15mm, breakable, opacityback=0, title=\textcolor{quarto-callout-important-color}{\faExclamation}\hspace{0.5em}{Qué exige (enunciado 2.2 y ficha del caso)}, toprule=.15mm, bottomtitle=1mm, colframe=quarto-callout-important-color-frame, rightrule=.15mm, titlerule=0mm, colback=white, arc=.35mm, left=2mm, opacitybacktitle=0.6, coltitle=black, leftrule=.75mm, colbacktitle=quarto-callout-important-color!10!white]

\begin{itemize}
\tightlist
\item
  Clave de fila y duplicados; filas y unidades antes y después de cada
  paso.
\item
  Auditar \texttt{inspection\_id}, licencia, dirección y cambios de
  nombre para enlazar historia \textbf{sin fusionar establecimientos
  distintos}.
\item
  Construir la respuesta: códigos, exclusiones; no asignar la categoría
  negativa a una respuesta desconocida.
\item
  Excluir \texttt{violations} de la inspección actual y cualquier texto
  que describa sus hallazgos.
\item
  Revisar calidad: faltantes, códigos especiales, categorías poco
  frecuentes, valores imposibles. Justificar toda exclusión.
\item
  Incluir una \textbf{tabla de construcción de la muestra} y un
  diagrama.
\end{itemize}

\end{tcolorbox}

\begin{longtable}[]{@{}lll@{}}

\caption{\label{tbl-construccion}Construcción de la muestra (versión
preliminar).}

\tabularnewline

\toprule\noalign{}
& Paso & Filas \\
\midrule\noalign{}
\endhead
\bottomrule\noalign{}
\endlastfoot
0 & Inspecciones en el paquete & 104.468 \\
1 & Con licencia válida & 104.216 \\
2 & Inspecciones Canvass (filas objetivo) & 52.675 \\
3 & Canvass con licencia válida & 52.548 \\

\end{longtable}

Se identificaron 252 inspecciones sin licencia válida (licencia
\texttt{0} o vacía) y 23.598 licencias válidas distintas. De las
inspecciones del paquete, 52.675 son de tipo \emph{Canvass}, cifra que
coincide con la ficha del caso.

\begin{tcolorbox}[enhanced jigsaw, toptitle=1mm, bottomrule=.15mm, breakable, opacityback=0, title=\textcolor{quarto-callout-note-color}{\faInfo}\hspace{0.5em}{Notas del equipo}, toprule=.15mm, bottomtitle=1mm, colframe=quarto-callout-note-color-frame, rightrule=.15mm, titlerule=0mm, colback=white, arc=.35mm, left=2mm, opacitybacktitle=0.6, coltitle=black, leftrule=.75mm, colbacktitle=quarto-callout-note-color!10!white]

\begin{itemize}
\tightlist
\item
  Falta agregar a la tabla: fusión de licencias sucesivas (23.418
  establecimientos) y construcción de la respuesta.
\item
  Decidir qué hacer con las 127 Canvass sin licencia válida.
\item
  Cuando exista \texttt{analisis.py}, esta tabla debe leerse de sus
  resultados y no recalcularse aquí.
\end{itemize}

\end{tcolorbox}

\subsection{Variables admisibles y
excluidas}\label{variables-admisibles-y-excluidas}

\begin{tcolorbox}[enhanced jigsaw, toptitle=1mm, bottomrule=.15mm, breakable, opacityback=0, title=\textcolor{quarto-callout-important-color}{\faExclamation}\hspace{0.5em}{Qué exige (enunciado 2.2.4)}, toprule=.15mm, bottomtitle=1mm, colframe=quarto-callout-important-color-frame, rightrule=.15mm, titlerule=0mm, colback=white, arc=.35mm, left=2mm, opacitybacktitle=0.6, coltitle=black, leftrule=.75mm, colbacktitle=quarto-callout-important-color!10!white]

Tabla con nombre, interpretación, tipo, disponibilidad temporal y
decisión de inclusión. Distinguir identificadores, datos de entrada,
resultados posteriores y variables derivadas de la respuesta. El nivel
de riesgo solo se usa si se justifica que antecede a la visita.

\end{tcolorbox}

\begin{tcolorbox}[enhanced jigsaw, toptitle=1mm, bottomrule=.15mm, breakable, opacityback=0, title=\textcolor{quarto-callout-note-color}{\faInfo}\hspace{0.5em}{Notas del equipo}, toprule=.15mm, bottomtitle=1mm, colframe=quarto-callout-note-color-frame, rightrule=.15mm, titlerule=0mm, colback=white, arc=.35mm, left=2mm, opacitybacktitle=0.6, coltitle=black, leftrule=.75mm, colbacktitle=quarto-callout-note-color!10!white]

\begin{itemize}
\tightlist
\item
\end{itemize}

\end{tcolorbox}

\subsection{Descripción de la
muestra}\label{descripciuxf3n-de-la-muestra}

\begin{tcolorbox}[enhanced jigsaw, toptitle=1mm, bottomrule=.15mm, breakable, opacityback=0, title=\textcolor{quarto-callout-important-color}{\faExclamation}\hspace{0.5em}{Qué exige (enunciado 2.2.6)}, toprule=.15mm, bottomtitle=1mm, colframe=quarto-callout-important-color-frame, rightrule=.15mm, titlerule=0mm, colback=white, arc=.35mm, left=2mm, opacitybacktitle=0.6, coltitle=black, leftrule=.75mm, colbacktitle=quarto-callout-important-color!10!white]

Distribución de la respuesta (solo desarrollo), cobertura temporal y
geográfica, número de entidades, frecuencia de eventos y pérdidas de
observaciones. Precisar a quiénes representan los datos.

\end{tcolorbox}

\begin{tcolorbox}[enhanced jigsaw, toptitle=1mm, bottomrule=.15mm, breakable, opacityback=0, title=\textcolor{quarto-callout-note-color}{\faInfo}\hspace{0.5em}{Notas del equipo}, toprule=.15mm, bottomtitle=1mm, colframe=quarto-callout-note-color-frame, rightrule=.15mm, titlerule=0mm, colback=white, arc=.35mm, left=2mm, opacitybacktitle=0.6, coltitle=black, leftrule=.75mm, colbacktitle=quarto-callout-note-color!10!white]

\begin{itemize}
\tightlist
\item
\end{itemize}

\end{tcolorbox}

\section{Diseño de evaluación}\label{sec-diseno}

\begin{tcolorbox}[enhanced jigsaw, toptitle=1mm, bottomrule=.15mm, breakable, opacityback=0, title=\textcolor{quarto-callout-important-color}{\faExclamation}\hspace{0.5em}{Qué exige (enunciado 4.1--4.2 y ficha del caso)}, toprule=.15mm, bottomtitle=1mm, colframe=quarto-callout-important-color-frame, rightrule=.15mm, titlerule=0mm, colback=white, arc=.35mm, left=2mm, opacitybacktitle=0.6, coltitle=black, leftrule=.75mm, colbacktitle=quarto-callout-important-color!10!white]

\begin{itemize}
\tightlist
\item
  Desarrollo 2019--2022 con ventanas temporales; reserva de prueba
  2023--2024, evaluada \textbf{una sola vez}.
\item
  Referencia: probabilidad constante de \emph{Fail} calculada en
  entrenamiento.
\item
  Métrica principal: pérdida logarítmica, más calibración.
  Complementarias: AUC, matriz de confusión, sensibilidad y VPP para un
  umbral fijado en desarrollo, prevalencias y tamaños.
\item
  Incluir un diagrama temporal de las particiones.
\end{itemize}

\end{tcolorbox}

\begin{tcolorbox}[enhanced jigsaw, toptitle=1mm, bottomrule=.15mm, breakable, opacityback=0, title=\textcolor{quarto-callout-note-color}{\faInfo}\hspace{0.5em}{Notas del equipo}, toprule=.15mm, bottomtitle=1mm, colframe=quarto-callout-note-color-frame, rightrule=.15mm, titlerule=0mm, colback=white, arc=.35mm, left=2mm, opacitybacktitle=0.6, coltitle=black, leftrule=.75mm, colbacktitle=quarto-callout-note-color!10!white]

\begin{itemize}
\tightlist
\item
\end{itemize}

\end{tcolorbox}

\section{Métodos}\label{sec-metodos}

\begin{tcolorbox}[enhanced jigsaw, toptitle=1mm, bottomrule=.15mm, breakable, opacityback=0, title=\textcolor{quarto-callout-important-color}{\faExclamation}\hspace{0.5em}{Qué exige (enunciado 3.1--3.2)}, toprule=.15mm, bottomtitle=1mm, colframe=quarto-callout-important-color-frame, rightrule=.15mm, titlerule=0mm, colback=white, arc=.35mm, left=2mm, opacitybacktitle=0.6, coltitle=black, leftrule=.75mm, colbacktitle=quarto-callout-important-color!10!white]

\begin{itemize}
\tightlist
\item
  M1: regresión logística. M2: árbol podado o SVM. M3: bosque aleatorio
  o boosting.
\item
  Expectativa previa sobre cada modelo: qué estructura podría capturar y
  cuándo podría fallar.
\item
  Preprocesamiento ajustado solo con datos de entrenamiento de cada
  partición.
\item
  Hiperparámetros: hasta 20 configuraciones por familia y tres
  particiones internas, salvo justificación.
\end{itemize}

\end{tcolorbox}

\begin{tcolorbox}[enhanced jigsaw, toptitle=1mm, bottomrule=.15mm, breakable, opacityback=0, title=\textcolor{quarto-callout-note-color}{\faInfo}\hspace{0.5em}{Notas del equipo}, toprule=.15mm, bottomtitle=1mm, colframe=quarto-callout-note-color-frame, rightrule=.15mm, titlerule=0mm, colback=white, arc=.35mm, left=2mm, opacitybacktitle=0.6, coltitle=black, leftrule=.75mm, colbacktitle=quarto-callout-note-color!10!white]

\begin{itemize}
\tightlist
\item
\end{itemize}

\end{tcolorbox}

\section{Resultados}\label{sec-resultados}

\begin{tcolorbox}[enhanced jigsaw, toptitle=1mm, bottomrule=.15mm, breakable, opacityback=0, title=\textcolor{quarto-callout-important-color}{\faExclamation}\hspace{0.5em}{Qué exige (enunciado 3.3, 4.2--4.4 y ficha del caso)}, toprule=.15mm, bottomtitle=1mm, colframe=quarto-callout-important-color-frame, rightrule=.15mm, titlerule=0mm, colback=white, arc=.35mm, left=2mm, opacitybacktitle=0.6, coltitle=black, leftrule=.75mm, colbacktitle=quarto-callout-important-color!10!white]

\begin{itemize}
\tightlist
\item
  Desarrollo y prueba claramente separados; tabla comparativa con
  diferencia frente a la referencia.
\item
  Incertidumbre que reconozca inspecciones repetidas del mismo
  establecimiento (remuestreo por establecimiento).
\item
  Errores por tipo de local y por historia disponible (con y sin
  antecedentes); una figura de diagnóstico.
\item
  Dos sensibilidades, por ejemplo: excluir antecedentes de inspección;
  evaluar \emph{Pass w/ Conditions} por separado; locales nuevos no
  vistos en entrenamiento.
\end{itemize}

\end{tcolorbox}

\begin{tcolorbox}[enhanced jigsaw, toptitle=1mm, bottomrule=.15mm, breakable, opacityback=0, title=\textcolor{quarto-callout-note-color}{\faInfo}\hspace{0.5em}{Notas del equipo}, toprule=.15mm, bottomtitle=1mm, colframe=quarto-callout-note-color-frame, rightrule=.15mm, titlerule=0mm, colback=white, arc=.35mm, left=2mm, opacitybacktitle=0.6, coltitle=black, leftrule=.75mm, colbacktitle=quarto-callout-note-color!10!white]

\begin{itemize}
\tightlist
\item
\end{itemize}

\end{tcolorbox}

\section{Discusión}\label{sec-discusion}

\begin{tcolorbox}[enhanced jigsaw, toptitle=1mm, bottomrule=.15mm, breakable, opacityback=0, title=\textcolor{quarto-callout-important-color}{\faExclamation}\hspace{0.5em}{Qué exige}, toprule=.15mm, bottomtitle=1mm, colframe=quarto-callout-important-color-frame, rightrule=.15mm, titlerule=0mm, colback=white, arc=.35mm, left=2mm, opacitybacktitle=0.6, coltitle=black, leftrule=.75mm, colbacktitle=quarto-callout-important-color!10!white]

Respuesta a la pregunta inicial, límites de generalización (solo locales
inspeccionados), utilidad y condiciones para usar el modelo. Una
situación en que la recomendación no sea confiable.

\end{tcolorbox}

\begin{tcolorbox}[enhanced jigsaw, toptitle=1mm, bottomrule=.15mm, breakable, opacityback=0, title=\textcolor{quarto-callout-note-color}{\faInfo}\hspace{0.5em}{Notas del equipo}, toprule=.15mm, bottomtitle=1mm, colframe=quarto-callout-note-color-frame, rightrule=.15mm, titlerule=0mm, colback=white, arc=.35mm, left=2mm, opacitybacktitle=0.6, coltitle=black, leftrule=.75mm, colbacktitle=quarto-callout-note-color!10!white]

\begin{itemize}
\tightlist
\item
\end{itemize}

\end{tcolorbox}

\section{Reproducibilidad y contribuciones}\label{sec-reproducibilidad}

\begin{tcolorbox}[enhanced jigsaw, toptitle=1mm, bottomrule=.15mm, breakable, opacityback=0, title=\textcolor{quarto-callout-important-color}{\faExclamation}\hspace{0.5em}{Qué exige (enunciado 5 y 6.4)}, toprule=.15mm, bottomtitle=1mm, colframe=quarto-callout-important-color-frame, rightrule=.15mm, titlerule=0mm, colback=white, arc=.35mm, left=2mm, opacitybacktitle=0.6, coltitle=black, leftrule=.75mm, colbacktitle=quarto-callout-important-color!10!white]

Comando de ejecución, versión del repositorio, entradas, tiempo de
ejecución, síntesis de aportes, declaración de uso de IA (herramienta,
fecha, propósito, partes afectadas y comprobaciones) y referencias.

\end{tcolorbox}

\begin{tcolorbox}[enhanced jigsaw, toptitle=1mm, bottomrule=.15mm, breakable, opacityback=0, title=\textcolor{quarto-callout-note-color}{\faInfo}\hspace{0.5em}{Notas del equipo}, toprule=.15mm, bottomtitle=1mm, colframe=quarto-callout-note-color-frame, rightrule=.15mm, titlerule=0mm, colback=white, arc=.35mm, left=2mm, opacitybacktitle=0.6, coltitle=black, leftrule=.75mm, colbacktitle=quarto-callout-note-color!10!white]

\begin{itemize}
\tightlist
\item
\end{itemize}

\end{tcolorbox}




\end{document}
